% Draft for paper 1's journal version (arXiv v2), hardening plan Phase 4, 11 Nov 2026. Not in the arXiv v1 source.
% Numbers: experiments/results/relay_noise/rn_metrics.tex (evaluation/relay_noise_macros.py from
% evaluation/relay_noise_sweep.py). To include: \input the macro file in the preamble and this paragraph after
% "Relaying rarely helps" in the swarm subsection.
\textbf{Noisy relays.} The relay hop above is clean apart from its static shadowing. We therefore swept
three kinds of noise:
\begin{itemize}
\item shadowing of 3, 6 and 8~dB;
\item per-packet loss of 0, 1, 5 and 10\,\% on every air-to-air hop, where a lost 1500-byte packet is resent until it
arrives, at the cost of its airtime plus a 1~ms backoff;
\item one or two air-to-air hops, the second through the intermediate drone with the best bottleneck rate, forwarded
store-and-forward.
\end{itemize}
We ran 5, 10 and 20 drones under both channel profiles, \RnSeeds{} seeds each, \RnCells{} relay settings in all. In none
of them did the order of the fixed split and the two adaptive policies change (\RnRankChanges{} changes). Frames
relayed by the adaptive policy stayed on time between \RnShareMin\,\% and \RnShareMax\,\% of the time, and its lead over the fixed split stayed between
\RnLeadMin{} and \RnLeadMax{} on-time frames/s per drone. Relaying itself still changed the adaptive policy's throughput
by only \RnGainMin{} to \RnGainMax{} frames/s (at most \RnGainRelPct\,\% of its rate), and a second hop by \RnSecondHopMin{} to \RnSecondHopMax{} (median \RnSecondHopMedian{}), so noise does not change the conclusion: the shared cell, not the relay
link, is the bottleneck.
